% !PRINTED 567
\weberpage{567}
%<<CONT

und nach dem binomischen Satze, da $a_n - b$ nicht negativ ist,
\[ \left(\frac{a_1 + a_2 + \cdots + a_n}{n}\right)^n \geqq b^n + b^{n-1}(a_n - b) = a_n b^{n-1}, \]
dies aber f\"allt mit der zu beweisenden Ungleichung (1) zusammen.

Es sei nun $\omega_1, \omega_2, \ldots, \omega_n$ eine Minimalbasis eines algebraischen
 K\"orpers $\Omega$. Es sei ferner $\varphi$ ein beliebiges ganzes Functional
.

Wir nehmen eine Basis von $\varphi$ an (\S.\,146):
\weq{2}{(\beta_1, \beta_2, \ldots, \beta_n) = \begin{pmatrix} b_{1,1} & \ldots & b_{1,n} \\ \cdot & \cdots & \cdot \\ b_{n,1} & \ldots & b_{n,n} \end{pmatrix} (\omega_1, \omega_2, \ldots, \omega_n),}
so dass
\weq{3}{N_a(\varphi) = \pm \Sigma \pm b_{1,1} \ldots b_{n,n}}
wird. Aus der mit $\varphi$ associirten Linearform
\weq{4}{\eta = \beta_1 x_1 + \beta_2 x_2 + \cdots + \beta_n x_n}
erhalten wir dann alle durch $\varphi$ theilbaren ganzen Zahlen des
K\"orpers $\Omega$, wenn f\"ur die $x_1, x_2, \ldots, x_n$ alle m\"oglichen Combinationen
 ganzer Zahlen gesetzt werden.

Aus (4) ergeben sich $n$ verschiedene Formen $\eta_1, \eta_2, \ldots, \eta_n$,
wenn wir zu den $n$ conjugirten K\"orpern, $\Omega_1, \Omega_2, \ldots, \Omega_n$ \"ubergehen
. Das Quadrat der Determinante dieses Systems von
$n$ Formen ist nach \S.\,147 gleich der Discriminante
\weq{5}{\Delta(\beta_1, \beta_2, \ldots, \beta_n) = \pm N_a(\varphi)^2 \Delta,}
wenn $\Delta$ die Grundzahl des K\"orpers bedeutet.

Die conjugirten K\"orper $\Omega_1, \Omega_2 \ldots$ k\"onnen theils reell, theils
paarweise conjugirt imagin\"ar sein. Wir bezeichnen, wie fr\"uher
schon mehrfach, mit $|\omega|$ den \gesperrt{absoluten Werth} einer reellen
oder complexen Gr\"osse und betrachten, indem wir unter
$x_1, x_2, \ldots, x_n$ reelle Variable verstehen, die quadratische Form
\weq{6}{f(x_1, x_2, \ldots, x_n) = |\eta_1|^2 + |\eta_2|^2 + \cdots + |\eta_n|^2 = S\,|\eta|^2.}
Ist $\Omega$ reell, so ist
\weq{7}{|\eta|^2 = \eta^2,}
also das Quadrat einer reellen linearen Form. Wenn aber $\Omega, \Omega'$
zwei conjugirt imagin\"are unter den conjugirten K\"orpern sind, so
sind auch die entsprechenden Linearformen $\eta, \eta'$ conjugirt imagin
\"ar, und es ist
\weq{8}{|\eta|^2 = |\eta'|^2 = \eta \eta'}
%CONT>>
